% Generated draft: requires mathematical and source review.
\documentclass[11pt]{article}
\usepackage[margin=25mm]{geometry}
\usepackage{fontspec}
\setmainfont{Latin Modern Roman}
\usepackage{amsmath,amssymb,mathtools,unicode-math}
\setmathfont{Latin Modern Math}
\usepackage{enumitem,needspace,xurl,hyperref}
\setlength{\parindent}{0pt}
\setlength{\parskip}{5pt}
\setlength{\emergencystretch}{3em}
\newcommand{\N}{\mathbb{N}}
\newcommand{\Z}{\mathbb{Z}}
\newcommand{\Q}{\mathbb{Q}}
\newcommand{\R}{\mathbb{R}}
\newcommand{\C}{\mathbb{C}}
\newcommand{\F}{\mathbb{F}}
\newcommand{\A}{\mathbb{A}}
\newcommand{\PP}{\mathbb{P}}
\newcommand{\E}{\mathbb{E}}
\newcommand{\eps}{\varepsilon}
\newcommand{\dd}{\,\mathrm{d}}
\begin{document}
\section*{UnsolvedMath ID: 30004674}
\textbf{Source problem code:} catalogue entry
\subsection{Definitions and mathematical statement}
Let \(\mathfrak c=2^{\aleph_0}=|\mathcal P(\N)|\). A partial order \((P,\leq_P)\) is a set equipped with a reflexive, antisymmetric, transitive relation. Interpret \emph{locally countable} to mean that every principal lower set is at most countable:
\[
\forall p\in P,\qquad |\{q\in P:q\leq_P p\}|\leq\aleph_0.
\]
The supplied statement does not explicitly define this term; this lower-set convention requires confirmation against the source.

For \(A,B\subseteq\N\), write \(A\leq_T B\) if an oracle Turing machine with oracle \(B\) computes the characteristic function of \(A\). Write \(A\equiv_T B\) if both \(A\leq_T B\) and \(B\leq_T A\). The set of Turing degrees is
\[
\mathcal D_T=\mathcal P(\N)/{\equiv_T},
\]
ordered by \([A]_T\leq_T[B]_T\) if \(A\leq_T B\). An order embedding \(f:P\to\mathcal D_T\) is an injective map satisfying
\[
\forall p,q\in P,\qquad p\leq_P q\ \Longleftrightarrow\ f(p)\leq_T f(q).
\]

\textbf{Question (attributed to Sacks in the supplied curation).} Does the following assertion hold in ZFC?
\[
\begin{gathered}
\text{For every partial order }(P,\leq_P)\text{ with }|P|=\mathfrak c
\text{ and }\forall p\in P\;
|\{q\in P:q\leq_P p\}|\leq\aleph_0,\\
\text{there exists an order embedding }f:(P,\leq_P)\longrightarrow(\mathcal D_T,\leq_T).
\end{gathered}
\]
No height restriction or additional set-theoretic axiom is imposed in this question. The supplied dated assessment labels the question open in ZFC and reports positive results under the continuum hypothesis, under Martin's axiom, and for restricted heights. These are supplied literature qualifications, not independently verified conclusions; the record does not specify the precise hypotheses or height bounds of those positive results.
\subsection{Short English statement}
Can every partial order of cardinality continuum, in which each element has at most countably many predecessors, be represented by Turing degrees so that the representation preserves and reflects the order? The interpretation of local countability uses lower sets. The supplied curation attributes the question to Sacks and labels it open; this statement-only draft requires human mathematical and source review.
\subsection{Sources}
\begin{itemize}
\item Supplied local UnsolvedMath registry record, permanent ID \(30004674\), title \emph{Embedding Locally Countable Orders into Turing Degrees}, exact source problem code \texttt{catalogue entry}, in the Oberwolfach Reports Open Problems collection. The record labels the statement \texttt{exact} and the problem \texttt{open}; these are supplied registry labels.
\item Supplied report citation: \emph{Computability Theory (2022)}, Oberwolfach Reports, DOI \texttt{10.4171/owr/2021/21}, pp.~1149--1190. \url{https://doi.org/10.4171/owr/2021/21}. The record gives proposed year 2021 and an \texttt{owr-year:2021} tag, while its citation uses 2022; this dating discrepancy is unresolved. No individual problem page or report passage is supplied.
\item Kumar, \emph{Set theory and Turing degrees}. Supplied literature reference: \url{https://home.iitk.ac.in/~krashu/tdst.pdf}.
\item Higuchi and Lutz, \emph{A note on a conjecture of Sacks}. Supplied literature reference: \url{https://arxiv.org/abs/2309.01876}.
\item Literature assessment dated \texttt{2026-08-22}: supplied dated curation, not independently verified. It attributes the question to Sacks, labels it open in ZFC, and reports positive results under CH, MA, and for restricted heights. The background describes this assessment as dated literature triage derived from the supplied OWR curation and requests independent verification of time-sensitive status and cited claims. The registry's \texttt{proposed\_by} field is empty, so the attribution here rests on that assessment.
\end{itemize}
\end{document}
